\subsection{The Last Days: 17 to 24 December}\label{sec:last-days}

On 17 December Advent changes pace. The Universal Norms on the Liturgical Year
order the weekdays from 17 to 24 December more directly to the preparation for
the Lord's birth (n.~42), and they rank them above the obligatory memorials
(n.~16 b); on them a saint's memorial may only be commemorated (General
Instruction of the Roman Missal, n.~355). So St Peter
Canisius on 21 December and St John of Kanty on 23 December are only
commemorated. From now on the Mass and the Office have their own texts for each
date, the priest sings the second Preface of Advent, and the readings, no longer
counted by the week, follow the calendar to the manger.

At Evening Prayer on each of the seven evenings from 17 to 23 December the
Church sings before the Canticle of Mary one of the great O antiphons. Each
calls on Christ by a name from the Old Testament and ends with ``come'': on 17
December \latin{O Sapientia}, O Wisdom; on the 18th \latin{O Adonai}, O Lord; on the
19th \latin{O radix Jesse}, O Root of Jesse; on the 20th \latin{O clavis David},
O Key of David; on the 21st \latin{O Oriens}, O Day-spring; on the 22nd \latin{O Rex
gentium}, O King of the nations; and on the 23rd \latin{O Emmanuel}. At Mass the
verse before the Gospel uses the same antiphons, shortened, and in a different
order, so that on the 21st the people hear
Emmanuel, on the 22nd and 23rd the King of the nations, and on the morning of
the 24th the Dawn from on high. The antiphons are given in full, in Latin and
English, in Appendix~\ref{sec:prayers}.

\subsubsection*{Day by day}

\textbf{17 December.} The first reading is the blessing of Jacob on Judah: ``The
sceptre shall not be taken away from Juda, nor a ruler from his thigh, till he
come that is to be sent, and he shall be the expectation of nations.'' The Gospel
is the genealogy of Christ at the head of St Matthew, the long roll of fathers
from Abraham to Joseph. In the Office of Readings Pope St Leo the Great writes
of the Incarnation:

\begin{quote}
And the fulfilment of the mystery of our atonement, which was ordained from all
eternity, was not assisted by any figures because the Holy Spirit had not yet
come upon the Virgin, and the power of the Most High had not over-shadowed her:
so that ``Wisdom building herself a house'' within her undefiled body, ``the
Word became flesh;'' and the form of God and the form of a slave coming together
into one person, the Creator of times was born in time; and He Himself through
whom all things were made, was brought forth in the midst of all things.\par
{\smallskip\noindent\small ---St Leo the Great, Letter 31\par}
\end{quote}

\noindent That evening the first of the great antiphons calls on the same
Wisdom, \latin{O Sapientia}, who comes ``out of the mouth of the Most High.''

\textbf{18 December.} Jeremiah promises that the Lord will ``raise up to David a
just branch,'' and the Gospel tells how the angel came to Joseph in his sleep:
``Joseph, son of David, fear not to take unto thee Mary thy wife, for that which
is conceived in her, is of the Holy Ghost.'' The Collect of the day grew from an
old Gelasian prayer that asked that we,
\latin{qui sub peccati iugo ex debito depremimur expectata Unigeniti tui nova
nativitate liberemur}, ``who are held down, as we deserve, under the yoke of sin,
may be set free by the long-awaited new birth of thy Only-begotten.'' The Office
reads the \emph{Letter to Diognetus}, an early defence of the faith by an
unknown Christian, whose great cry of wonder is given in
Section~\ref{sec:exchange}.

\textbf{19 December.} The angel promises a son to the barren wife of Manoah, and
in the Gospel another angel promises a son to old Zachary in the Temple. In the
Office St Irenaeus speaks of the sign Isaiah promised:

\begin{quote}
On this account, therefore, the Lord Himself, who is Emmanuel from the Virgin,
is the sign of our salvation, since it was the Lord Himself who saved them,
because they could not be saved by their own instrumentality\par
{\smallskip\noindent\small ---St Irenaeus, \emph{Against Heresies}, III.20\par}
\end{quote}

\noindent That evening the antiphon calls on the Root of Jesse,
\latin{qui stas in signum populorum}, who stands as a sign for the peoples. Both
Irenaeus and the antiphon draw on Isaiah, who promised a sign in the Virgin and
saw the root of Jesse stand ``for an ensign of the people.''

\textbf{20 December.} The prophecy of the Virgin who shall conceive is read with
the Gospel of the Annunciation. Whether the 20th is a weekday or the Fourth
Sunday, the Office of Readings reads St Bernard on the moment when heaven waited
for Mary's reply:

\begin{quote}
Make haste, therefore, to answer the Angel, or rather to answer the Lord through
the Angel. Say the word and receive the Word. Utter thy human word and conceive
the Divine Word. Pronounce the transitory word and embrace the Word everlasting.
Why dost thou hesitate? Wherefore dost thou fear? \ldots{} Behold the Desired of
all nations is standing outside and knocking at thy door. \ldots{} Arise,
therefore, and make haste to open to Him. Arise by faith, make haste by
devotion, open by consent.\par
{\smallskip\noindent\small ---St Bernard, Fourth Homily on the \latin{Missus est}\par}
\end{quote}

\noindent Bernard's Latin gives the command in three words, \latin{Surge,
curre, aperi}, ``Arise, run, open.''

\textbf{21 December.} The first reading may be from the Song of Songs or from
Zephaniah, and the Gospel is the Visitation. Setting the Song beside Mary's
journey is the Lectionary's own work:

\begin{quote}
The voice of my beloved, behold he cometh leaping upon the mountains, skipping
over the hills.\par
{\smallskip\noindent\small ---Song of Songs 2:8 (the Canticle of Canticles in the Douay)\par}
\end{quote}

\noindent In the Office St Ambrose, as the \emph{Catena Aurea} gives him, asks
why Mary went in haste: ``For what could Mary now, filled with God, but ascend
into the higher parts with haste!'' And he turns her song to us: ``But let every
one have the spirit of Mary, so that he may rejoice in the Lord.''

\textbf{22 December.} Hannah brings the child Samuel to the Lord at Shiloh:
``For this child did I pray, and the Lord hath granted me my petition, which I
asked of him.'' Her song is the psalm, and the Gospel is the Magnificat. In the
Office St Bede the Venerable, as the \emph{Catena Aurea} gives him, explains why Mary's
spirit rejoices: ``Because the
spirit of the Virgin rejoices in the eternal Godhead of the same Jesus, (i. e.
the Saviour,) whose flesh is formed in the womb by a temporal conception.''

\textbf{23 December.} Malachi foretells the messenger, ``Behold I send my angel,
and he shall prepare the way before my face,'' and the return of Elijah before
the great day of the Lord (Malachi 3:23 in the Lectionary, 4:5 in the Douay).
The Gospel tells of John's birth and of Zachary's tongue loosed when he wrote,
``John is his name.'' In the Office St Hippolytus confesses the Trinity revealed
by the Incarnation: ``We accordingly see the Word incarnate, and we know the
Father by Him, and we believe in the Son, (and) we worship the Holy Spirit.''

\textbf{24 December.} In the Office of Readings of the last morning, St
Augustine begins a Christmas sermon with a sentence that returns to the Wisdom
first invoked a week before:
\latin{Natalis Domini dicitur, quando Dei Sapientia se demonstravit infantem, et
Dei Verbum sine verbis vocem carnis emisit}, ``It is called the Lord's birthday
when the Wisdom of God showed himself an infant, and the Word of God, without
words, uttered the voice of flesh.'' A little further on he cries,
\latin{Expergiscere, homo: pro te Deus factus est homo}, ``Awake, O man: for thee
God was made man.''

The Collect of that morning is the last cry of Advent. It reworks a prayer of
the Gelasian Advent Masses, which the Gregorian book, in the copy called the
Hadrianum, set on the Ember Wednesday of December, with the station at St Mary
Major. The Gelasian form begins:
\latin{Festina, quaesumus, ne tardaveris, Domine}, ``Hasten, we beseech thee, do
not delay, O Lord.'' The present Collect turns the same plea to the Lord Jesus
himself.
